% =====================================================================
% §2.4 — La deviazione per grano: ampiezza e probabilità
%
% Esempi:
%   deviation  -> (a)/(b) (stream mask_range + mask_probability) in UN
%                      solo YAML; UN solo \lstinputlisting che copre i due
%                      stream (linerange da verificare sul file reale);
%                      figura a due pannelli deviation_map.
%   probability     -> quarto angolo: deviation_probability globale [[0,0],[1,100]] SENZA
%                      alcun range -> jitter implicito (micromodulazione).
%                      Mostrato inline (una sola chiave); figura opzionale.
%   exA_micro       -> controparte UDIBILE (freeze dry/wet) nel bundle.
%
% Pattern: aggancio all'identità di contenuto di §2.3 -> (a)/(b) (la sola
%   posizione dell'envelope, offset_range come istanza) -> lettura figura
%   -> meccanismo: due equazioni a confronto (Truax eq:truax vs gated
%   eq:gated) -> quadrato 2x2 (envelope puro / Truax / gated / quarto
%   angolo = jitter implicito) con tabella dei jitter -> controparte sul
%   contenuto della distribution di §2.3 -> riproducibilità -> ponte §2.5
%
% Dipendenze di preambolo:
%   - amsmath  (equation, \eqref, \tfrac)
%   - textcomp (\textdegree)  [già usato altrove]
%   - capt-of/caption (\captionof)  [già usato per le figure]
%   - tabella: \hline (niente booktabs); se booktabs è disponibile,
%     \toprule/\midrule/\bottomrule rendono meglio.
%
% TODO prima della submission:
%   [ ] linerange del listato unico: verificare i numeri sul file reale
%       (deve coprire entrambi gli stream, mask_range + mask_probability)
%   [ ] \ref del rinvio di generalità (parametri con range+deviation_probability):
%       allineare al label reale della sezione che li dispiega (sec:voci?)
%   [x] tabella jitter: ora generata da examples/gen_jitter_table.py dai
%       default_jitter del PGE pinnato (make jitter-table). Verificato:
%       volume/pan/durata/lettura = ±default_jitter/2 (UniformDistribution);
%       pitch ±12 cents = EDO_IMPLICIT_DETUNE_CENTS (semi-ampiezza diretta);
%       footnote ratio corretta a ±0,0025 (era ±0,005, errata di un fattore 2).
%   [ ] \cite{Truax1988}, \cite{Vaggione2002}: chiavi reali
%   [ ] \ref{fig:pointer-MAP}, \ref{sec:griglia}, \ref{sec:pointer},
%       \ref{sec:voci}: verificare label
%   [ ] figura deviation_map a colonna singola: leggibilità a 8,2 cm
%       (due pannelli impilati); valutare figure* se troppo piccola
% =====================================================================

% ----------------------------------------------------------------
\subsection{La deviazione per grano: ampiezza e probabilità}\label{sec:deviazione}

Quasi ogni parametro del sistema è una funzione del tempo dichiarata in
due parti: una \emph{base} (la traiettoria centrale) e un \emph{range} di
deviazione attorno ad essa, campionato in modo indipendente per ogni
grano. Si può inoltre decidere la \emph{probabilità} con cui il range si applichi grano per grano; se non
dichiarata è aperta al 100\%, ogni grano devia. Le tre chiavi hanno lo
stesso nome per ogni parametro che le ammette --- posizione di lettura
(\texttt{offset\_range}), volume, pan, durata, pitch --- e gli esempi
successivi le dispiegano su ciascuno (\S\ref{sec:voci}). Questa
sottosezione isola le due chiavi indipendenti, ampiezza del range e
probabilità, sulla sola posizione di lettura.

I due stream di Fig.~\ref{fig:deviazione-ab} differiscono per la sola posizione
dell'inviluppo. In (a)
(\texttt{mask\_range}) l'inviluppo sta sul \emph{range}: cresce da zero al
35\% del buffer, con \texttt{deviation\_probability} assente e dunque
sempre attiva. In (b) (\texttt{mask\_probability}) il range resta fisso al
massimo di (a) e a crescere da 0 a 100\% è la
\texttt{deviation\_probability}.

\begin{figure}[H]
    \begin{framedfloat}
    \lstinputlisting[
      linerange={44,45,50-53,55,57,58,63-67},
      language=yaml,
      basicstyle=\ttfamily\fontsize{7}{8.4}\selectfont
    ]{examples/deviation/deviation.yml}
    \vspace{2pt}
    \includegraphics[width=\linewidth]{examples/deviation/deviation_annotated}
    \caption{In alto, i due stream in un solo YAML (\texttt{deviation.yml}):
    base condivisa (\texttt{freeze} a 0.5, \texttt{speed\_ratio} nullo) e
    stessa banda massima; \texttt{mask\_range} pone l'inviluppo
    sull'ampiezza (\texttt{offset\_range}), \texttt{mask\_probability} sulla
    probabilità (\texttt{deviation\_probability.pointer}).
    Sotto, le due maschere della deviazione in un'unica \textsc{map}
    (assi come in
    Fig.~\ref{fig:c-e}). Stessa base, stessa banda massima: ampiezza
    e probabilità sono assi indipendenti. Pannello (a)
    (\texttt{mask\_range}): il range si allarga con continuità e i grani lo
    riempiono uniformemente, un cuneo in cui tutti deviano, di poco
    all'inizio e di molto alla fine. Pannello (b)
    (\texttt{mask\_probability}): il range resta al massimo fin dall'inizio,
    ma solo una popolazione crescente di grani lo usa --- non un
    allargamento ma una mistura fra grani fedeli alla linea centrale e
    grani devianti, in proporzione che cresce nel tempo.}
    \label{fig:deviazione-ab}
    \end{framedfloat}
\end{figure}

Il meccanismo sottostante è la \emph{tendency mask} della tradizione
granulare~\cite{Truax1988}. Per ogni parametro essa assegna al grano $n$,
di onset normalizzato $\tau_n$, una traiettoria centrale $c$ e un range di
deviazione $\rho$, entrambi inviluppi, con un campionamento
indipendente per grano:
\begin{equation}
  v_n \;=\; c(\tau_n) \;+\; \xi_n\,\rho(\tau_n),
  \qquad \xi_n \sim \mathcal{U}\!\left(-\tfrac{1}{2},\,\tfrac{1}{2}\right),
  \label{eq:truax}
\end{equation}
dove $v_n$ è il valore che il parametro assume per il grano $n$ e $\xi_n$
il campione estratto per quel grano. Il caso (a) è questo modello
all'opera, con la traiettoria centrale $c(\tau_n)$ costante --- la lettura
è ferma --- e il solo range $\rho$ che varia nel tempo. Il contributo
proprio del sistema è un secondo
fattore, anch'esso componibile come inviluppo: un gate di probabilità $g_n$
che decide grano per grano \emph{se} la maschera si applica,
\begin{equation}
  v_n \;=\; c(\tau_n) \;+\; g_n\,\xi_n\,\rho(\tau_n),
  \qquad g_n \sim \mathrm{Bernoulli}\!\left(p(\tau_n)\right),
  \label{eq:gated}
\end{equation}
con $p$ la chiave \texttt{deviation\_probability}. La sola differenza fra le due è il
fattore $g_n$: per $p\equiv 1$ si ha $g_n\equiv 1$ e la eq.~\eqref{eq:gated}
ricade nella eq.~\eqref{eq:truax}. Ampiezza ($\rho$) e probabilità ($p$) sono
dunque due inviluppi indipendenti, e i due stream muovono un inviluppo per
volta tenendo fermo l'altro: in (a) cresce $\rho$ mentre $p$ resta a 1
(il gate è sempre aperto, ogni grano devia); in (b) cresce $p$ mentre
$\rho$ resta al valore massimo (la banda è fissa, cambia solo quanti
grani la usano).

La \texttt{deviation\_probability} funziona però anche senza alcun range
esplicito (Fig.~\ref{fig:probability}): dichiarata da sola, apre un
jitter implicito di entità minima, fissata dal sistema, su ciascuno dei
parametri elencati in Tab.~\ref{tab:jitter}
--- una micromodulazione che non disegna alcuna traiettoria ma decorrela
la massa grano per grano, nei termini della \emph{décorrélation
microtemporelle} di Vaggione~\cite{Vaggione2002}.

\begin{figure}[H]
    \begin{framedfloat}
    \lstinputlisting[
      linerange={19-22,26},
      language=yaml,
      basicstyle=\ttfamily\fontsize{7}{8.4}\selectfont
    ]{examples/probability/probability.yml}
    \vspace{2pt}
    \includegraphics[width=\linewidth]{examples/probability/probability_map.pdf}
    \caption{In alto, la sola \texttt{deviation\_probability} aggiunta alle
    condizioni minime, senza alcun range (\texttt{probability.yml}).
    Sotto, la \textsc{map}
    (stream \texttt{probability\_ramp}; assi come in
    Fig.~\ref{fig:c-e}). La probabilità sale da~0 a~100\% a gradini,
    in quattro scatti fra il 20\% e l'80\% della durata; la banda netta di
    sinistra si sfrangia in nuvola a destra.}
    \label{fig:probability}
    \end{framedfloat}
\end{figure}

\begin{table}[H]
  \centering
  \footnotesize
  \input{figures/jitter_table}
  \normalsize
  \caption{Jitter implicito per parametro quando la
  \texttt{deviation\_probability} è dichiarata senza range esplicito.}
  \label{tab:jitter}
\end{table}

